% Generated draft: requires mathematical and source review.
\documentclass[11pt]{article}
\usepackage[margin=25mm]{geometry}
\usepackage{fontspec}
\setmainfont{Latin Modern Roman}
\usepackage{amsmath,amssymb,mathtools,unicode-math}
\setmathfont{Latin Modern Math}
\usepackage{enumitem,needspace,xurl,hyperref}
\setlength{\parindent}{0pt}
\setlength{\parskip}{5pt}
\setlength{\emergencystretch}{3em}
\newcommand{\N}{\mathbb{N}}
\newcommand{\Z}{\mathbb{Z}}
\newcommand{\Q}{\mathbb{Q}}
\newcommand{\R}{\mathbb{R}}
\newcommand{\C}{\mathbb{C}}
\newcommand{\F}{\mathbb{F}}
\newcommand{\A}{\mathbb{A}}
\newcommand{\PP}{\mathbb{P}}
\newcommand{\E}{\mathbb{E}}
\newcommand{\eps}{\varepsilon}
\newcommand{\dd}{\,\mathrm{d}}
\begin{document}
\section*{UnsolvedMath ID: 30005389}
\textbf{Source problem code:} catalogue entry
\subsection{Definitions and mathematical statement}
The supplied record identifies the following assertion as Conjecture 1. Its ambient Shimura-variety and Hecke-theoretic setup is not specified completely, so the formulation below is conditional on that intended setup.

Let \(X_K(\C)\) denote the complex analytic space at level \(K\) occurring in the source, and let \(d\) denote its designated middle cohomological degree. For a smooth, pure-dimensional complex variety, this would be its complex dimension; the supplied passage does not explicitly define \(d\). Let \(\ell\) be the coefficient prime and \(\F_\ell\) the field with \(\ell\) elements. Write \(H^i(X_K(\C),\F_\ell)\) for ordinary singular cohomology. Let \(\mathbb T\) denote the intended commutative Hecke algebra acting on these groups, and let \(\mathfrak m\subset\mathbb T\) be the intended maximal ideal. For a \(\mathbb T\)-module \(M\), write
\[
M_{\mathfrak m}=(\mathbb T\setminus\mathfrak m)^{-1}M.
\]
The choice of \(\mathbb T\), including its coefficient ring and the places at which Hecke operators are included, is not supplied.

Assume that the residual Galois representation \(\overline\rho_{\mathfrak m}\) attached to \(\mathfrak m\) is defined in the intended setup. Absolute irreducibility means that it remains irreducible after extending its coefficient field to an algebraic closure. The record uses \emph{non-Eisenstein} to mean precisely this absolute irreducibility. Its Galois group, coefficient field, dimension, and existence hypotheses are not specified in the excerpt.

\textbf{Conjectural target.} For every level and maximal ideal to which the intended source setup applies,
\[
\overline\rho_{\mathfrak m}\text{ absolutely irreducible}
\quad\Longrightarrow\quad
H^i(X_K(\C),\F_\ell)_{\mathfrak m}=0
\qquad\text{for every integer }i\geq 0\text{ with }i\neq d.
\]
This concentration assertion does not require nonvanishing in degree \(d\). The supplied record does not determine the full class of varieties, levels, or coefficient primes over which this implication is conjectured.

\textbf{Additional source context.} The passage describes this as a special case of a general folklore conjecture about cohomology of locally symmetric spaces, with potential applications to the Langlands program. It then discusses understanding the structure of the mod-\(\ell\) groups through a generalization of Ihara's lemma, with a conjectural description at least away from \(\ell\). In that subsequent context, it takes a prime \(\mathfrak p\mid p\) of a field \(F\), assumes \(D_{\mathfrak p}\) is split and
\[
K=K^{\mathfrak p}K_{\mathfrak p},
\qquad K_{\mathfrak p}=\operatorname{GL}_2(\mathcal O_{F,\mathfrak p}),
\]
and introduces the subgroup with Iwahori level
\[
K_0(\mathfrak p)=
\left\{(k^{\mathfrak p},k_{\mathfrak p})\in K:
 k_{\mathfrak p}\equiv
 \begin{pmatrix} *&0\\ *&*\end{pmatrix}
 \pmod{\mathfrak p}\right\}.
\]
It gives two degeneracy maps \(\pi_1,\pi_2:X_{K_0(\mathfrak p)}\to X_K\), used for the Hecke correspondence for \(T_{\mathfrak p}\). The record does not define \(F\), \(D\), or the local ring notation fully, nor does it state a further Ihara-type conjecture. These subsequent local hypotheses are not explicitly imposed on Conjecture 1, and the excerpt does not explicitly require \(p\neq\ell\).
\subsection{Short English statement}
The conjecture predicts that ordinary mod-\(\ell\) cohomology localized at a Hecke maximal ideal vanishes outside the designated middle degree whenever the attached residual Galois representation is absolutely irreducible. The supplied excerpt omits the ambient geometric and Hecke hypotheses needed to specify its full scope. Its supplied dated assessment labels the problem partially solved; this draft does not independently verify that assessment.
\subsection{Sources}
\begin{itemize}
\item Local UnsolvedMath registry, permanent ID \(30005389\), source problem code \texttt{catalogue entry}, title \emph{Middle-Degree Concentration of Non-Eisenstein Cohomology}. The record belongs to the Oberwolfach Reports Open Problems collection and gives proposed year 2023; no individual proposer is supplied.
\item Supplied primary-source citation: \emph{Arithmetic of Shimura Varieties} (2023), Oberwolfach Reports, DOI 10.4171/owr/2023/5, pp.\ 261--326, \url{https://doi.org/10.4171/owr/2023/5}. The excerpt labels the target Conjecture 1. No precise page within the report or individual contribution attribution is supplied.
\item Supplied literature source: Caraiani and Tamiozzo, \emph{Hilbert modular varieties with torsion coefficients}, \url{https://arxiv.org/abs/2107.10081}.
\item Supplied literature source: \emph{Oberwolfach Report 5/2023}, \url{https://doi.org/10.4171/owr/2023/5}.
\item Literature assessment dated exactly \texttt{2026-08-22}: the supplied curation labels the problem partially solved and reports middle-degree concentration in important generic Hilbert-modular and related cases, including non-solvable residual image hypotheses, while describing the general conjecture as open. This is supplied dated curation, not an independently verified literature assessment. The background expressly requests independent verification of time-sensitive status and cited claims.
\item The registry supplies the statement-status label \texttt{corrected\_verified} and says the statement was reconstructed and checked against the official Oberwolfach report. These are supplied provenance claims; no source consultation or independent mathematical verification is asserted here.
\end{itemize}
\end{document}
